% DP305 user manual -- example for the plaindoc class.
% Orrin Instruments and the DP305 are fictional. Build: latexmk -lualatex dp305-manual.tex
\documentclass[type=manual, lang=en-GB, otherlangs={zh-HK}]{plaindoc}

\pdset{
  brand    = {Orrin Instruments},
  title    = {DP305 Bench Power Supply},
  subtitle = {User manual},
  product  = {DP305},
  doctype  = {User manual},
  docid    = {DP305-UM-EN},
  version  = {Version 1.2},
  date     = {25 November 2026},
  org      = {Orrin Instruments is a fictional company. The DP305 and its
              specifications exist only to demonstrate the plaindoc class.},
}
% Technical drawings for the DP305 example manual.
% Coordinates in millimetres. Line roles come from the pd/* TikZ styles.

% Front panel. #1 = scale, #2 = 1 to draw callouts.
\newcommand\dpfront[2][1]{%
\begin{tikzpicture}[pd/drawing, x=#1mm, y=#1mm]
  % feet and case
  \foreach \x in {14,122} \draw[pd/detail, fill=pd-surface] (\x,-2.2) rectangle ++(14,2.2);
  \draw[pd/part, rounded corners=3mm] (0,0) rectangle (150,66);
  \draw[pd/detail, rounded corners=1.8mm] (3,3) rectangle (147,63);
  % display window
  \draw[pd/detail, fill=pd-surface, rounded corners=1.2mm] (8,28) rectangle (80,58);
  \node[anchor=base east, font=\sffamily\fontsize{21}{21}\selectfont] at (58,45) {12.00};
  \node[anchor=base west, font=\sffamily\fontsize{11}{11}\selectfont] at (59.5,45) {V};
  \node[anchor=base east, font=\sffamily\fontsize{21}{21}\selectfont] at (58,32) {0.250};
  \node[anchor=base west, font=\sffamily\fontsize{11}{11}\selectfont] at (59.5,32) {A};
  \node[fill=pd-ink, text=white, inner sep=1pt, minimum width=7.5mm, font=\sffamily\fontsize{5.5}{5.5}\selectfont\bfseries] at (71.5,52.5) {CV};
  \node[draw=pd-ink2, text=pd-ink2, inner sep=1pt, minimum width=7.5mm, line width=.3pt, font=\sffamily\fontsize{5.5}{5.5}\selectfont\bfseries] at (71.5,47) {CC};
  \node[fill=pd-ink, text=white, inner sep=1pt, minimum width=7.5mm, font=\sffamily\fontsize{5.5}{5.5}\selectfont\bfseries] at (71.5,41.5) {OUT};
  \node[draw=pd-ink2, text=pd-ink2, inner sep=1pt, minimum width=7.5mm, line width=.3pt, font=\sffamily\fontsize{5.5}{5.5}\selectfont\bfseries] at (71.5,36) {LOCK};
  % power switch
  \draw[pd/detail, rounded corners=.8mm] (9,8) rectangle (23,20);
  \draw[pd/detail, fill=pd-surface, rounded corners=.5mm] (11,10) rectangle (21,18);
  \node[pd/label, font=\sffamily\fontsize{5}{5}\selectfont] at (16,23) {POWER};
  % keys
  \foreach \x/\t in {32/LIMIT, 49/LOCK, 66/OUTPUT}{
    \draw[pd/detail, rounded corners=.8mm] (\x,9) rectangle ++(14,8);
    \node[font=\sffamily\fontsize{5}{5}\selectfont\bfseries] at (\x+7,13) {\t};}
  \fill[pd-ink] (73,20.5) circle (.9);
  % knobs
  \foreach \cx/\t in {98/VOLTAGE, 126/CURRENT}{
    \node[pd/label, font=\sffamily\fontsize{5}{5}\selectfont] at (\cx,30.2) {\t};
    \draw[pd/detail] (\cx,43) circle (10);
    \draw[pd/part] (\cx,43) circle (8);
    \draw[pd/detail] (\cx,43) circle (6.2);
    \draw[pd/detail, line width=.8pt] (\cx,46.5) -- (\cx,50);}
  % output terminals
  \foreach \tx/\t in {100/+, 117/−, 134/GND}{
    \draw[pd/part] (\tx,15) circle (4.4);
    \draw[pd/detail, fill=pd-surface] (\tx,15) circle (2.1);
    \node[pd/label, font=\sffamily\fontsize{6}{6}\selectfont\bfseries] at (\tx,23) {\t};}
  \ifnum#2=1
    \pdcallout{1}{(30,58)}{(30,76)}
    \pdcallout{2}{(98,53)}{(98,76)}
    \pdcallout{3}{(126,53)}{(126,76)}
    \pdcallout{4}{(39,9)}{(39,-11)}
    \pdcallout{5}{(56,9)}{(56,-11)}
    \pdcallout{6}{(73,9)}{(73,-11)}
    \draw[pd/leader] (100,10.6) -- (100,7) -- (117,7) -- (117,10.6);
    \pdcallout{7}{(108.5,7)}{(108.5,-11)}
    \pdcallout{8}{(134,10.6)}{(134,-11)}
    \pdcallout{9}{(16,8)}{(16,-11)}
  \fi
\end{tikzpicture}}

% Rear panel with callouts.
\newcommand\dprear{%
\begin{tikzpicture}[pd/drawing, x=1mm, y=1mm]
  \foreach \x in {14,122} \draw[pd/detail, fill=pd-surface] (\x,-2.2) rectangle ++(14,2.2);
  \draw[pd/part, rounded corners=3mm] (0,0) rectangle (150,66);
  \draw[pd/detail, rounded corners=1.8mm] (3,3) rectangle (147,63);
  % fan grille
  \draw[pd/detail] (38,33) circle (21);
  \foreach \r in {16,11,6} \draw[pd/detail] (38,33) circle (\r);
  \draw[pd/detail] (17,33) -- (59,33) (38,12) -- (38,54);
  \fill[pd-ink] (38,33) circle (2.2);
  % rating plate
  \draw[pd/detail] (70,44) rectangle (98,58);
  \foreach \y in {54.5,51.5,48.5} \draw[pd-rule, line width=1.1pt] (73,\y) -- (95,\y);
  % USB-C
  \draw[pd/part, rounded corners=1.4mm] (118,50) rectangle ++(10,3.6);
  \node[pd/label, font=\sffamily\fontsize{5}{5}\selectfont] at (123,57.5) {USB};
  % AC inlet (IEC 60320 C14)
  \draw[pd/part, rounded corners=1mm] (106,20) rectangle (132,40);
  \draw[pd/detail, fill=pd-surface] (110,24) -- (128,24) -- (128,33) -- (125,36) -- (113,36) -- (110,33) -- cycle;
  \foreach \px/\py in {113.5/28.5, 124.5/28.5, 119/32.5} \fill[pd-ink] (\px-1,\py-1.6) rectangle ++(2,3.2);
  % fuse drawer
  \draw[pd/detail, rounded corners=.6mm] (110,9) rectangle (128,16);
  \draw[pd/detail] (116,12.5) -- (122,12.5);
  \node[pd/label, font=\sffamily\fontsize{4.6}{4.6}\selectfont] at (119,5.6) {T3.15AL 250V};
  \pdcallout{1}{(38,54)}{(38,76)}
  \pdcallout{2}{(84,58)}{(84,76)}
  \pdcallout{3}{(123,53.6)}{(123,76)}
  \pdcallout{4}{(132,30)}{(158,30)}
  \pdcallout{5}{(128,12.5)}{(158,12.5)}
\end{tikzpicture}}

% Connection: supply terminals to a load, current direction shown.
\newcommand\dpconnect{%
\begin{tikzpicture}[pd/drawing, x=1mm, y=1mm]
  % supply (terminal strip)
  \draw[pd/part, rounded corners=1.5mm] (0,0) rectangle (52,30);
  \node[pd/label, anchor=north west] at (3,28) {DP305};
  \foreach \tx/\t in {12/+, 28/−, 44/GND}{
    \draw[pd/part] (\tx,11) circle (4.4);
    \draw[pd/detail, fill=pd-surface] (\tx,11) circle (2.1);
    \node[pd/label, font=\sffamily\fontsize{6}{6}\selectfont\bfseries] at (\tx,20) {\t};}
  % load
  \draw[pd/part, rounded corners=1.5mm] (100,-4) rectangle (140,34);
  \node[pd/label] at (120,28) {Load};
  \draw[pd/detail] (120,4) -- (120,8);
  \draw[pd/part, fill=white] (116,8) rectangle (124,22);
  \draw[pd/detail] (120,22) -- (120,26);
  \node[font=\sffamily\fontsize{4.6}{4.6}\selectfont, anchor=west] at (125,15) {max.\ 5\,A};
  % wires
  \draw[line width=1.3pt, pd-ink] (12,11) -- (12,-12) -- (90,-12) -- (90,26) -- (120,26);
  \draw[line width=1.3pt, pd-ink, dash pattern=on 4pt off 1.6pt] (28,11) -- (28,-20) -- (95,-20) -- (95,4) -- (120,4);
  \foreach \p in {(12,11),(28,11),(120,26),(120,4)} \fill[pd-ink] \p circle (.9);
  \node[font=\sffamily\fontsize{5.5}{5.5}\selectfont, anchor=south] at (60,-12) {positive lead};
  \node[font=\sffamily\fontsize{5.5}{5.5}\selectfont, anchor=north] at (60,-20) {negative lead};
  \draw[pd/dim, ->] (64,-7.5) -- (80,-7.5);
  \node[font=\sffamily\fontsize{5}{5}\selectfont, text=pd-ink2, anchor=west] at (81,-7.5) {I};
\end{tikzpicture}}

\pdcover{\centering\dpfront[1.05]{0}}

\begin{document}
\maketitle
\tableofcontents

% ------------------------------------------------------------------
\chapter{About this manual}

This manual describes how to set up and operate the DP305, a single-output
laboratory power supply rated \qtyrange{0}{30}{\volt} and \qtyrange{0}{5}{\ampere}.
It is written for people who connect the supply to electronic circuits and
assumes basic knowledge of voltage, current and polarity.

\section{Conventions}

\begin{tabularx}{\linewidth}{@{}P{30mm}L@{}}
\toprule
\hd{Notation} & \hd{Meaning} \\
\midrule
\key{OUTPUT}   & A key on the front panel. \\
\ui{VOLTAGE}   & A label printed on the instrument or shown on the display. \\
\menu{System>Interface} & A menu path: open \ui{System}, then select \ui{Interface}. \\
\cmd{MEAS:VOLT?} & A remote command, typed exactly as shown. \\
→ & The result you should see after a step. \\
\bottomrule
\end{tabularx}

\section{Related documents}

The \emph{DP305 Programming Reference} (DP305-PR-EN) lists every remote
command. The \emph{DP305 Service Manual} (DP305-SM-EN) is for qualified
service personnel only.

% ------------------------------------------------------------------
\chapter{Safety}

Read this chapter before you connect the DP305 to mains power or to a load.
Safety messages in later chapters repeat the hazards that apply to a
specific step.

\section{Safety messages}

Each safety message starts with a signal word. The signal word tells you how
serious the hazard is.

\begin{tabularx}{\linewidth}{@{}P{30mm}L@{}}
\toprule
\hd{Signal word} & \hd{Meaning} \\
\midrule
DANGER  & A hazardous situation that \emph{will} cause death or serious injury if not avoided. \\
WARNING & A hazardous situation that \emph{could} cause death or serious injury if not avoided. \\
CAUTION & A hazardous situation that could cause minor or moderate injury if not avoided. \\
NOTICE  & A situation that could damage the instrument, the load or data. No personal injury. \\
\bottomrule
\end{tabularx}

\section{Before you use the DP305}

\begin{safety}{warning}[Electric shock from damaged mains connections]
A damaged power cord or inlet can expose mains voltage.
\begin{itemize}
\item Use only the supplied power cord, or one rated for at least
      \qty{250}{\volt} and \qty{10}{\ampere} with a protective earth conductor.
\item Do not use the DP305 if the cord, inlet or case is damaged.
\end{itemize}
\end{safety}

\begin{safety}{caution}[Burns from wires and loads]
At \qty{5}{\ampere}, thin wires and some loads become hot enough to burn skin.
Use wires with a cross-section of at least \qty{1.0}{\mm\squared}.
Switch the output off and let the load cool before you touch it.
\end{safety}

\begin{safety}{notice}
Blocking the fan or the side vents makes the DP305 switch its output off
through over-temperature protection. Keep \qty{10}{\cm} of free space
behind the instrument.
\end{safety}

\section{Intended use}

The DP305 supplies DC power to electronic circuits in laboratories, schools
and workshops. Do not use it to charge batteries, to power medical devices
connected to a person, or in wet or explosive environments.

% ------------------------------------------------------------------
\chapter{Product overview}

\section{Front panel}

\begin{figure}[h]
\pdwide{\centering\dpfront{1}}
\caption{Front panel. Numbers refer to the list below.}
\label{fig:front}
\end{figure}

\begin{pdlegend}
\legenditem{1}{Display}{Output voltage and current, mode and status indicators.}
\legenditem{2}{VOLTAGE knob}{Sets the output voltage in \qty{10}{\mV} steps. Press to switch to \qty{1}{\volt} steps.}
\legenditem{3}{CURRENT knob}{Sets the current limit while \key{LIMIT} is active.}
\legenditem{4}{LIMIT key}{Shows the set values instead of the measured values.}
\legenditem{5}{LOCK key}{Locks the knobs and keys. Hold for \qty{2}{\s} to unlock.}
\legenditem{6}{OUTPUT key}{Switches the output on or off. The lamp above it lights while the output is on.}
\legenditem{7}{Output terminals}{Positive (+) and negative (−) terminals. Accept \qty{4}{\mm} plugs, spade lugs and bare wire.}
\legenditem{8}{GND terminal}{Connected to protective earth. Not connected to the output.}
\legenditem{9}{POWER switch}{Switches the instrument on or off. The output always starts off.}
\end{pdlegend}

\section{Display indicators}

\begin{tabularx}{\linewidth}{@{}P{16mm}L@{}}
\toprule
\hd{Indicator} & \hd{Meaning} \\
\midrule
\textbf{CV}   & Constant voltage. The output voltage equals the set value. \\
\textbf{CC}   & Constant current. The load draws the current limit, so the voltage is lower than set. \\
\textbf{OUT}  & The output is on. \\
\textbf{LOCK} & The front panel is locked. \\
\textbf{OTP}  & Over-temperature protection switched the output off. See \cref{ch:trouble}. \\
\bottomrule
\end{tabularx}

\section{Rear panel}

\begin{figure}[h]
\pdwide{\centering\dprear}
\caption{Rear panel.}
\label{fig:rear}
\end{figure}

\begin{pdlegend}
\legenditem{1}{Fan outlet}{Keep clear. The fan runs faster as the load increases.}
\legenditem{2}{Rating plate}{Serial number and mains rating.}
\legenditem{3}{USB-C port}{Remote control. Does not power the DP305.}
\legenditem{4}{Mains inlet}{IEC 60320 C14 inlet for the power cord.}
\legenditem{5}{Fuse drawer}{Holds the mains fuse, T3.15AL 250V, and one spare.}
\end{pdlegend}

% ------------------------------------------------------------------
\chapter{Setting up}

\section{Before you begin}

You need:
\begin{itemize}
\item a mains socket with protective earth, \qtyrange{100}{240}{\volt} AC, \qty{50}{\Hz} or \qty{60}{\Hz};
\item test leads rated for at least \qty{5}{\ampere};
\item the voltage and current your load needs, from its documentation.
\end{itemize}

\section{Switch on}

\begin{steps}
\item Place the DP305 on a flat surface with \qty{10}{\cm} of free space behind it.
\item Connect the power cord to the mains inlet and to an earthed socket.
\item Press the \ui{POWER} switch.
      \result{The display shows the last set values. The \ui{OUT} indicator is off.}
\end{steps}

\section{Connect a load}
\label{sec:connect}

\begin{figure}[h]
\pdwide{\centering\dpconnect}
\caption{Connect the positive lead first and disconnect it last.}
\label{fig:connect}
\end{figure}

\begin{safety}{caution}[Reversed polarity]
Many loads are damaged when + and − are swapped. Check the polarity marking
on the load before you connect it.
\end{safety}

\begin{steps}
\item If \ui{OUT} is lit, press \key{OUTPUT} to switch the output off.
\item Set the voltage and current limit before connecting (see \cref{sec:set}).
\item Connect the positive lead from the + terminal to the positive input of the load.
\item Connect the negative lead from the − terminal to the negative input of the load.
\item Press \key{OUTPUT}.
      \result{\ui{OUT} lights. The display shows the measured voltage and current.}
\end{steps}

\begin{note}
The GND terminal is not part of the output circuit. Connect it to the load
only if the load's documentation asks for an earth connection.
\end{note}

% ------------------------------------------------------------------
\chapter{Operation}

\section{Set voltage and current limit}
\label{sec:set}

\begin{steps}
\item Press \key{LIMIT}.
      \result{The display shows the set values. \ui{CV} and \ui{CC} go off.}
\item Turn the \ui{VOLTAGE} knob to the voltage your load needs.
\item Turn the \ui{CURRENT} knob to the highest current the load should draw.
\item Press \key{LIMIT} again, or wait \qty{5}{\s}.
      \result{The display returns to the measured values.}
\end{steps}

\begin{note}
Set the current limit just above the load's normal current. If a fault
occurs, the DP305 then limits the current before the load is damaged.
\end{note}

\section{Constant voltage and constant current}

The DP305 changes between the two modes by itself. While the load draws less
than the current limit, the supply holds the voltage (\ui{CV}). When the load
tries to draw more, the supply holds the current and lowers the voltage
(\ui{CC}).

\begin{tabularx}{\linewidth}{@{}P{24mm}LL@{}}
\toprule
\hd{Indicator} & \hd{What is happening} & \hd{What to do} \\
\midrule
\textbf{CV} & Normal operation. & Nothing. \\
\textbf{CC} & The load needs more current than the limit, or it is short-circuited. &
  Check the load. Raise the limit only if the load is rated for the higher current. \\
\bottomrule
\end{tabularx}

\section{Lock the front panel}

Press \key{LOCK}. \ui{LOCK} appears and the knobs and keys stop responding,
except \key{OUTPUT}, which can still switch the output off. Hold \key{LOCK}
for \qty{2}{\s} to unlock.

\section{Remote control over USB}

The DP305 appears as a virtual serial port. Commands follow SCPI syntax and
end with a line feed.

\begin{tabularx}{\linewidth}{@{}P{34mm}L@{}}
\toprule
\hd{Command} & \hd{Function} \\
\midrule
\cmd{*IDN?}          & Returns manufacturer, model, serial number and firmware version. \\
\cmd{VOLT <value>}   & Sets the output voltage in volts. \\
\cmd{CURR <value>}   & Sets the current limit in amperes. \\
\cmd{OUTP ON|OFF}    & Switches the output on or off. \\
\cmd{MEAS:VOLT?}     & Returns the measured output voltage. \\
\cmd{MEAS:CURR?}     & Returns the measured output current. \\
\bottomrule
\end{tabularx}

Example session (sent lines start with \cmd{>}):
\begin{code}
> VOLT 5.00
> CURR 0.200
> OUTP ON
> MEAS:CURR?
0.1482
\end{code}

% ------------------------------------------------------------------
\chapter{Troubleshooting}
\label{ch:trouble}

\begin{xltabular}{\linewidth}{@{}P{36mm}LL@{}}
\toprule
\hd{Symptom} & \hd{Possible cause} & \hd{Action} \\
\midrule
\endhead
Display stays dark. & No mains power, or the fuse has blown. &
  Check the socket and cord. Replace the fuse with the spare in the fuse drawer. \\
\addlinespace
\ui{CC} lights and the voltage is low. & The load draws the current limit. &
  Check the load for a short circuit. Raise the limit if the load needs it. \\
\addlinespace
\ui{OTP} lights and the output is off. & The DP305 is too hot. &
  Clear the fan outlet and vents. Wait \qty{10}{\minute}, then press \key{OUTPUT}. \\
\addlinespace
Knobs do nothing. & The front panel is locked. & Hold \key{LOCK} for \qty{2}{\s}. \\
\addlinespace
No serial port on the computer. & Charge-only USB cable, or no driver. &
  Use a USB data cable. The port uses the standard CDC-ACM driver. \\
\bottomrule
\end{xltabular}

If none of these actions helps, note the serial number from the rating plate
and the firmware version from \cmd{*IDN?}, then contact your supplier.

% ------------------------------------------------------------------
\chapter{Specifications}

Values apply after \qty{30}{\minute} of warm-up at \qtyrange{18}{28}{\degreeCelsius}.

\begin{xltabular}{\linewidth}{@{}LP{58mm}@{}}
\toprule
\hd{Parameter} & \hd{Value} \\
\midrule
\endhead
\pdgroup{Output}
Voltage                      & \qtyrange{0}{30}{\volt} \\
Current                      & \qtyrange{0}{5}{\ampere} \\
Setting resolution           & \qty{10}{\mV}, \qty{1}{\mA} \\
Voltage accuracy             & ±(\qty{0.05}{\percent} + \qty{10}{\mV}) \\
Current accuracy             & ±(\qty{0.1}{\percent} + \qty{5}{\mA}) \\
Ripple and noise (\qty{20}{\Hz} to \qty{20}{\MHz}) & ≤ \qty{1}{\mV} rms \\
Load regulation              & ≤ \qty{0.01}{\percent} + \qty{2}{\mV} \\
\addlinespace
\pdgroup{Mains supply}
Voltage                      & \qtyrange{100}{240}{\volt} AC, \qtylist{50;60}{\Hz} \\
Power consumption            & ≤ \qty{250}{\volt\ampere} \\
Fuse                         & T3.15AL 250V, \qtyproduct{5 x 20}{\mm} \\
\addlinespace
\pdgroup{General}
Interface                    & USB-C (USB 2.0, CDC-ACM), SCPI \\
Dimensions (W × H × D)       & \qtyproduct{215 x 90 x 260}{\mm} \\
Mass                         & \qty{3.1}{\kg} \\
Operating temperature        & \qtyrange{0}{40}{\degreeCelsius} \\
Storage temperature          & \qtyrange{-20}{60}{\degreeCelsius} \\
\bottomrule
\end{xltabular}

% ------------------------------------------------------------------
\begin{lang}{zh-HK}
\chapter{快速參考（繁體中文）}

本章摘要 DP305 的基本操作。開始前，請先閱讀第 2 章的安全資訊。

\begin{safety}{warning}[輸出端子帶有危險能量]
輸出端子最高可提供 \qty{30}{\volt}、\qty{5}{\ampere}。短路會令導線發熱，並可能導致灼傷。
連接或拆除負載前，請先按 \key{OUTPUT} 關閉輸出。
\end{safety}

\section{設定並開啟輸出}

\begin{steps}
\item 按下 \ui{POWER} 開關。
      \result{顯示屏顯示上次的設定值，輸出保持關閉。}
\item 按 \key{LIMIT}，然後轉動 \ui{VOLTAGE} 旋鈕設定電壓。
\item 轉動 \ui{CURRENT} 旋鈕，設定負載可承受的最大電流。
\item 將正極導線接到 + 端子，負極導線接到 − 端子。
\item 按 \key{OUTPUT}。
      \result{\ui{OUT} 指示燈亮起，顯示屏顯示實際電壓及電流。}
\end{steps}

\section{指示燈}

\begin{tabularx}{\linewidth}{@{}P{16mm}L@{}}
\toprule
\hd{指示燈} & \hd{意思} \\
\midrule
\textbf{CV}  & 定電壓：輸出電壓等於設定值。 \\
\textbf{CC}  & 定電流：負載電流已達上限，因此電壓低於設定值。請檢查負載是否短路。 \\
\textbf{OTP} & 過熱保護：輸出已關閉。清除通風口的阻擋物，等候 \qty{10}{\minute} 後再按 \key{OUTPUT}。 \\
\bottomrule
\end{tabularx}
\end{lang}

\end{document}
